% META: DONE.1

\section{Применение машинного обучения в CAD}

\subsection{Направления исследований на стыке ML и CAD}

За последние годы сформировалось несколько активно развивающихся направлений на стыке машинного обучения и систем автоматизированного проектирования:

\begin{itemize}
    \item Генерация CAD-моделей. Нейросетевые методы используются для автоматической генерации параметрических моделей по изображениям, эскизам или текстовому описанию. Это включает генерацию последовательности операций (например, выдавливаний) для восстановления 3D-формы.
    \item Классификация 3D-объектов. Распознавание типа детали (фланец, вал, корпус) по её геометрии - задача, для которой успешно применяются PointNet и её производные.
    \item Семантическая сегментация 3D-геометрии. Разделение модели на семантически значимые части (например, выделение отверстий, пазов, рёбер жёсткости) - прямая аналогия с задачей, решаемой в данной работе.
    \item Оптимизация топологии. Использование обучения с подкреплением для поиска оптимальной конструкции с заданными прочностными и весовыми характеристиками.
    \item Восстановление истории построения. Восстановление последовательности операций по итоговой геометрии - наиболее сложная задача, сочетающая элементы классификации, сегментации и последовательностного моделирования.
\end{itemize}

\subsection{Ключевые работы и полученные результаты}

Помимо уже упомянутых работ Autodesk Research, следует отметить несколько значимых исследований:

DeepCAD (2021) \cite{stanfordrlcad2021}. Исследователи из Стэнфорда предложили подход к генерации CAD-моделей как последовательностей операций. Нейросетевая модель обучалась на датасете из 178 000 моделей и могла генерировать реалистичные последовательности выдавливаний и вырезаний. Однако работа фокусировалась на генерации, а не на восстановлении истории по готовой геометрии.

CAD-SIGNet (2023) \cite{Khan_2024_CVPR}. Работа, посвящённая извлечению параметрических сигнатур из B-Rep представления с использованием графовых нейронных сетей. Авторы показали, что GNN могут эффективно распознавать типы поверхностей и их взаимосвязи 

Segmenting CAD models with MeshCNN (2022) \cite{10908022}. Ряд работ демонстрирует применение MeshCNN для сегментации CAD-моделей на семантические части, такие как отверстия, пазы, скругления. Эти исследования напрямую подтверждают применимость выбранной архитектуры для задачи, решаемой в данной работе.

\subsection{Проблемы и открытые вопросы}

Несмотря на впечатляющий прогресс, применение машинного обучения в CAD сталкивается с рядом проблем, которые остаются открытыми:

\begin{itemize}
    \item Нехватка размеченных данных. Создание датасетов CAD-моделей с полной историей построения требует значительных усилий. Даже Autodesk, имея доступ к миллионам моделей Fusion 360, публично выпустила лишь небольшую часть.
    \item Проблема обобщения. Модели, обученные на одних типах деталей, часто плохо работают на других, а обучение на простых геометриях не гарантирует успеха на сложных промышленных объектах.
    \item Чувствительность к численным параметрам. Как показано в данной работе, нейросети хорошо распознают тип операции, но испытывают трудности с точной регрессией численных параметров; это фундаментальное ограничение, которое требует гибридных подходов.
    \item Интерпретируемость. Для промышленного применения важно не только получить результат, но и понять, почему нейросеть приняла то или иное решение. CAD-инженеры должны доверять системе, что невозможно без прозрачности. Как минимум, система должна предоставлять возможность детализации и отображения отдельных промежуточных результатов, чтобы инженер имел возможность разобраться в том или ином ответе системы и научиться её понимать.
    \item Интеграция с существующими CAD-системами. Большинство исследований проводятся в изолированной среде (Python + библиотеки для 3D). Интеграция обученных моделей в промышленные CAD-системы, такие как КОМПАС-3D, требует дополнительных инженерных усилий.
\end{itemize}

